\subsection*{Binding two names to the roles requested on a pictured page}
A reader can return familiar words in grammatical order and still credit the wrong person. We exchanged two names in each of 112 held-out Chinese passages and asked about each role under each assignment. Both pages contain the same multiset of glyphs. Each credited answer must identify the work, relation and participant in a complete statement. A writer trained on genuine question--answer pairs satisfied all four requests for 40 passages. Exchanging questions on a random half of the training boards reduced this count to two. Rephrasing the evaluation questions reduced the genuine writer's count to five. Training on these pages therefore conveyed some dependence on the requested fact, with pronounced sensitivity to its wording.

Combining two individually practised facts exposed a larger failure (Fig.~\ref{fig:grounded-proposition}). When composing both sentences in a single call, the writer returned 118 of 312 complete replies for practised combinations of roles and three of 136 for withheld combinations. The latter were lyricist with composer and director with screenwriter. Varying the wording during additional training raised the latter count to eleven. A learned visual planner instead identified the headword and the two requested descriptions and passed each description to the same writer. This system returned 168 of 312 practised replies and 70 of 136 withheld replies. The paired difference from the single-call ordered writer on withheld combinations was 49.3 percentage points (95\% passage-bootstrap interval, 34.6--64.7). It satisfied both orders on both assignments for 15 of the 34 withheld passages. Its selected input spans were exact on 444 of the 448 compound requests. Supplying the annotated spans gave the same 70 withheld replies. Reversing the selected roles gave zero. Correctly identifying the requested roles thus accounts for only part of a successful answer.

The first held-out passage produced two correct credits for a song before and after the name exchange. A separately authored film note gave a different result. Although the planner returned a correctly formed sentence about each of the two roles, it assigned both to the director. After exchanging the names it also copied the verb preceding the name (Fig.~\ref{fig:grounded-proposition}d). On the 32 individual programme questions the planned writer produced 14 complete accepted statements. A descriptive reading of local text Qwen3-8B's answers found the requested fact in all 32, although none followed the enumerated forms exactly. On 29 original CMRC and 19 DRCD questions, the planned writer included the reference name in 20 and six responses, respectively, but produced only two and zero complete accepted statements. The restricted scorer can miss sound paraphrases. The programme images independently exhibit actual contradictions and attribution errors. On the original open evaluation the writer matched ten of 571 references and none of the 300 requiring content ink absent from the input.

We separated evidence selection from sentence construction by placing a reader-selected name ahead of the original passage. The visual span reader located 40 of the 48 human reference answers exactly. Given those predicted spans, a writer trained with genuine evidence produced three complete accepted statements. Supplying all 48 correct spans still yielded three. Some answers were useful direct quotations. Others assigned a person's name to both the work and its director. On the short programme notes the same intervention yielded 15 of 32 accepted statements with predicted evidence and 28 with correct evidence. Thus a misplaced span accounted for many programme errors, whereas supplying the correct span alone seldom recovered a complete proposition on the broader human questions.
